\begin{table}[htbp]
\centering
\caption{Summary Statistics by Procedure Type and Period}
\label{tab:summary}
\begin{threeparttable}
\begin{tabular}{lcccc}
\toprule
& \multicolumn{2}{c}{Ballot Cantons} & \multicolumn{2}{c}{Administrative Cantons} \\
\cmidrule(lr){2-3} \cmidrule(lr){4-5}
& Pre-ruling & Post-ruling & Pre-ruling & Post-ruling \\
\midrule
Naturalization rate & 18.1 & 20.8 & 21.2 & 24.6 \\
\quad (per 1,000 foreign) & (42.7) & (25.7) & (52.1) & (31.2) \\
Mean naturalizations & 9.5 & 22.3 & 12.1 & 26.2 \\
Mean foreign population & 654 & 1085 & 820 & 1188 \\
Mean total population & 3992 & 4776 & 3191 & 3755 \\
Foreign share (\%) & 10.4 & 16.1 & 14.6 & 19.1 \\
Observations & 20,369 & 18,823 & 12,054 & 11,147 \\
Municipalities & 900 & 900 & 531 & 531 \\
\bottomrule
\end{tabular}
\begin{tablenotes}
\item \textit{Notes:} Standard deviations in parentheses. Ballot cantons are German-speaking cantons (ZH, LU, SZ, OW, NW, GL, ZG, SO, BL, SH, AR, AI, SG, AG, TG) where municipalities used popular assembly votes for naturalization decisions before the 2003 Federal Court ruling. Administrative cantons (GE, VD, NE, TI, JU, BS) already used written administrative procedures. Pre-ruling: 1981--2003. Post-ruling: 2004--2024. Naturalization rate = annual naturalizations per 1,000 foreign residents.
\end{tablenotes}
\end{threeparttable}
\end{table}

